% ============================================================
%  BAB 2 — LANDASAN TEORI (REVISI LENGKAP)
%  NusaNara: Sistem Bimbingan Karier Berbasis RAG dan LLM
% ============================================================
\chapter{LANDASAN TEORI}

Bab ini menyajikan fondasi teoretis yang menjadi dasar bagi keseluruhan penelitian.
Uraian disusun secara sistematis, dimulai dari konsep perancangan dan validasi sistem,
menukik ke teknologi inti yang digunakan, membahas kerangka teoretis dari domain
bimbingan karir, hingga memetakan lanskap riset untuk memposisikan kebaruan penelitian ini.

% -------------------------------------------------------
\section{Perancangan Sistem}
% -------------------------------------------------------
Perancangan sistem merupakan tahapan krusial dalam \textit{System Development Life Cycle}
yang mentransformasikan kebutuhan fungsional dan non-fungsional menjadi struktur teknis
dan arsitektur sistem. Proses ini mencakup pemodelan data menggunakan
\textit{Entity--Relationship Diagram} dan Normalisasi, desain antarmuka pengguna yang
ergonomis, penentuan komponen perangkat keras serta perangkat lunak, serta integrasi modul
untuk memastikan kinerja, skalabilitas, keamanan, dan kemudahan pemeliharaan.
Tujuannya adalah menghasilkan spesifikasi teknis yang meminimalkan risiko proyek,
memenuhi kebutuhan pengguna dengan efisiensi biaya dan waktu, serta mendukung tujuan
strategis organisasi melalui integrasi sistem yang optimal \cite{Varezki2021}.

% -------------------------------------------------------
\section{Validasi Sistem}
% -------------------------------------------------------
Validasi adalah proses yang dilakukan untuk memastikan bahwa sistem atau perangkat lunak
yang dikembangkan benar-benar memenuhi dan menjawab kebutuhan pengguna.
Validasi perangkat lunak merupakan fase utama dalam rekayasa kebutuhan yang memastikan
bahwa kebutuhan sesuai dengan sistem target, dengan tujuan mendeteksi dan memperbaiki
kesalahan dalam kebutuhan yang telah ditentukan sebelumnya \cite{Atoum2021}.

% -------------------------------------------------------
\section{Design Science Research (DSR)}
% -------------------------------------------------------
Menurut \citet{Peffers2007}, \textit{Design Science Research}
(DSR) merupakan pendekatan penelitian yang berorientasi pada penciptaan dan evaluasi
artefak guna memecahkan masalah organisasi yang teridentifikasi melalui proses yang
sistematis dan \textit{rigor}. Proses tersebut meliputi: (1) \textit{problem identification},
(2) \textit{objectives of a solution}, (3) \textit{design \& development},
(4) \textit{demonstration}, (5) \textit{evaluation}, dan (6) \textit{communication}.

Output penelitian DSR berupa artefak dapat berwujud \textit{construct} (konsep baru),
\textit{model} (representasi), \textit{method} (prosedur atau algoritma), maupun
\textit{instantiation} (implementasi sistem).

\begin{figure}[H]
  \centering
  \missingfig{Proses Design Science Research (Peffers et al., 2007) -- Gambar 2.1 dari draf asli}
  \caption{Proses \textit{Design Science Research} (Sumber: Peffers et al., 2007)}
  \label{fig:dsr}
\end{figure}

% -------------------------------------------------------
\section{Sistem Adaptif}
% -------------------------------------------------------
Sistem adaptif merupakan sistem yang mampu menyesuaikan perilaku, konfigurasi,
atau strukturnya secara otomatis sebagai respons terhadap perubahan kondisi internal
maupun eksternal \cite{Saputri2020}. Sistem cerdas modern menggunakan prinsip adaptabilitas
ini untuk mempelajari pola interaksi secara dinamis dan memperbarui status agen (memori) tanpa memerlukan modifikasi sistem secara manual \cite{Russell2021}.

% -------------------------------------------------------
\section{Personalisasi}
% -------------------------------------------------------
Personalisasi merupakan bentuk adaptasi di mana sistem menyesuaikan konten, layanan,
atau antarmuka berdasarkan karakteristik, preferensi, dan kebutuhan spesifik pengguna \cite{Benfarha2023}.
Personalisasi tingkat lanjut (\textit{hyper-personalization}) mengintegrasikan jejak historis
dan profil semantik untuk menghasilkan output yang kontekstual, sehingga relevansi
informasi yang diberikan kepada setiap individu menjadi optimal \cite{Adomavicius2005}.

% -------------------------------------------------------
\section{Artificial Intelligence (AI)}
% -------------------------------------------------------
Kecerdasan buatan (\textit{Artificial Intelligence}, AI) adalah bidang ilmu komputer
yang mempelajari dan merealisasikan agen atau sistem yang mampu menunaikan
tugas-tugas yang biasanya membutuhkan kecerdasan manusia, termasuk pengenalan pola,
inferensi, pembuatan keputusan, dan pembelajaran dari data \cite{Alhejaily2025}.

% -------------------------------------------------------
\section{Large Language Models (LLM)}
% -------------------------------------------------------
\textit{Large Language Models} (LLM) adalah model bahasa berbasis jaringan saraf
berskala sangat besar yang dilatih pada korpus teks masif untuk memprediksi token
berikutnya dalam sebuah urutan \cite{Brown2020}. Peningkatan skala model secara signifikan
memperluas kemampuan generalisasi, termasuk \textit{few-shot learning}.

Studi menunjukkan bahwa performa LLM mengikuti \textit{scaling laws}, di mana kualitas
model meningkat secara konsisten mengikuti \textit{power-law} terhadap ukuran model,
jumlah data, dan komputasi pelatihan \cite{Kaplan2020}.

\subsection{Local LLM dan Platform Ollama}
Perkembangan teknik kuantisasi model (seperti format GGUF dan GPTQ) memungkinkan LLM berskala
miliaran parameter dijalankan secara lokal pada perangkat keras komoditas (\textit{commodity hardware}) tanpa
memerlukan infrastruktur klaster komputasi pusat. Paradigma pemrosesan lokal ini didukung oleh berbagai \textit{runtime engine}, salah satunya adalah
\textbf{Ollama}, yang menyederhanakan siklus hidup pengoperasian LLM terbuka melalui
antarmuka baris perintah dan layanan HTTP API \cite{Ollama2024}.

Keunggulan pendekatan lokal dibandingkan Cloud API disajikan pada Tabel \ref{tab:local_vs_cloud}:

\begin{table}[H]
\centering
\caption{Perbandingan Local LLM vs. Cloud API}
\label{tab:local_vs_cloud}
\begin{tabularx}{\textwidth}{lXX}
\toprule
\textbf{Aspek} & \textbf{Cloud API (Gemini, OpenAI)} & \textbf{Local LLM (Ollama)} \\
\midrule
Privasi data & Data dikirim ke server eksternal & Data tetap di infrastruktur lokal \\
Biaya operasional & \textit{Per-token billing} & Gratis setelah setup awal \\
Latensi & Bergantung jaringan & Bergantung hardware lokal \\
Kontrol & Terbatas, bergantung vendor & Penuh, termasuk parameter model \\
Ketersediaan offline & Tidak & Ya \\
\bottomrule
\end{tabularx}
\end{table}

% -------------------------------------------------------
\section{Retrieval-Augmented Generation (RAG)}
% -------------------------------------------------------
\textit{Retrieval-Augmented Generation} (RAG) adalah arsitektur yang menggabungkan
proses \textit{retrieval} (pengambilan) dokumen dari basis pengetahuan eksternal dengan
kemampuan model generatif untuk menghasilkan jawaban yang lebih akurat dan
kontekstual \cite{Lewis2020}.

Alur kerja RAG dimulai ketika pengguna mengajukan pertanyaan (1), yang memicu
pencarian semantik untuk mengambil dokumen relevan dari basis pengetahuan (2).
Dokumen tersebut digabungkan dengan \textit{prompt} asli untuk memperkaya konteks (3),
kemudian model memproses informasi tersebut untuk menghasilkan respons akhir (4).

\begin{figure}[H]
  \centering
  \missingfig{Arsitektur RAG -- Gambar 2.2 dari draf asli}
  \caption{Arsitektur Model \textit{Retrieval-Augmented Generation} (RAG) (Sumber: diadaptasi dari \cite{Lewis2020})}
  \label{fig:rag}
\end{figure}

Pendekatan ini terbukti efektif mengurangi \textit{hallucination} karena model tidak
hanya mengandalkan memori internal, tetapi juga referensi eksternal yang dinamis \cite{Lewis2020}.

% -------------------------------------------------------
\section{Hybrid Search dan Reciprocal Rank Fusion (RRF)}
% -------------------------------------------------------
\subsection{Konsep Hybrid Search}
\textit{Hybrid Search} menggabungkan dua paradigma pencarian yang saling komplementer:
pencarian semantik berbasis vektor (\textit{dense retrieval}) dan pencarian berbasis
kata kunci (\textit{sparse/lexical retrieval}). Masing-masing memiliki kelemahan
inheren \cite{Luan2021}:

\begin{itemize}
  \item Pencarian semantik: memahami sinonim dan konteks, tetapi lemah terhadap \textit{exact match} istilah teknis.
  \item Pencarian kata kunci: sangat tepat untuk istilah spesifik, tetapi tidak memahami variasi makna.
\end{itemize}

Dalam sistem manajemen basis data modern, keduanya dapat diimplementasikan secara terintegrasi;
misalnya, pencarian semantik menggunakan fungsi jarak spasial (\textit{Cosine Distance}),
dan pencarian kata kunci menggunakan indeks \textit{Generalized Inverted Index} (GIN) \cite{Wang2021}.

\subsection{Reciprocal Rank Fusion (RRF)}
\textit{Reciprocal Rank Fusion} (RRF) adalah algoritma \textit{rank aggregation} yang
menggabungkan beberapa \textit{ranked list} menjadi satu peringkat tunggal tanpa
normalisasi skor eksplisit \cite{Cormack2009}. Rumus RRF adalah:

\begin{equation}
  RRF\_score(d) = \frac{1}{k + rank_{semantic}(d)} + \frac{1}{k + rank_{lexical}(d)}
  \label{eq:rrf}
\end{equation}

\noindent di mana:
\begin{itemize}[noitemsep]
  \item $d$ = dokumen kandidat
  \item $rank_{semantic}(d)$ = peringkat dokumen $d$ dalam hasil pencarian semantik
  \item $rank_{lexical}(d)$ = peringkat dokumen $d$ dalam hasil pencarian kata kunci
  \item $k$ = konstanta \textit{smoothing} ($k = 60$, divalidasi oleh Cormack et al., 2009)
\end{itemize}

Dokumen yang muncul di peringkat tinggi pada \textbf{kedua} sistem mendapat skor RRF
signifikan lebih tinggi, menghasilkan daftar kandidat final yang tangguh terhadap kelemahan
masing-masing metode \textit{retrieval} secara individu.

% -------------------------------------------------------
\section{Server-Sent Events (SSE)}
% -------------------------------------------------------
\textit{Server-Sent Events} (SSE) adalah standar web API yang memungkinkan server
mengirimkan data ke klien secara \textit{unidirectional} melalui koneksi HTTP yang
persisten \cite{W3C2015}. SSE dirancang khusus untuk aliran data satu arah dari server
ke klien, menjadikannya pilihan yang lebih sederhana dan efisien dibandingkan WebSocket
untuk kasus penggunaan \textit{streaming} teks.

Dalam konteks LLM, SSE menjadi mekanisme standar untuk mengirimkan token respons
satu per satu secara \textit{real-time}, sehingga pengguna dapat mulai membaca
respons sebelum LLM selesai menghasilkan keseluruhan teks. Hal ini secara signifikan
menurunkan persepsi latensi (\textit{perceived latency}).

Secara teknis, aliran data SSE diformat menggunakan \texttt{text/event-stream}:

\begin{lstlisting}[language=bash, caption={Struktur Dasar Payload SSE}]
Content-Type: text/event-stream
Cache-Control: no-cache
Connection: keep-alive

data: Data chunk 1

data: Data chunk 2

data: [DONE]
\end{lstlisting}

% -------------------------------------------------------
\section{Model Embedding dan Representasi Vektor}
% -------------------------------------------------------
Model \textit{embedding} adalah pilar utama dalam pemrosesan bahasa alami modern yang
bertugas memetakan teks (kata, kalimat, atau dokumen) ke dalam ruang vektor
berdimensi tinggi, sehingga mesin dapat menghitung metrik kesamaan (\textit{similarity})
secara matematis \cite{Mikolov2013}.

Perkembangan arsitektur \textit{embedding} bergeser dari model \textit{dense} klasik
menuju \textit{Mixture-of-Experts} (MoE). Arsitektur MoE memungkinkan model memiliki
skala parameter yang besar namun membatasi jumlah komputasi aktif per token dengan merutekan
pemrosesan ke sub-jaringan (\textit{expert}) tertentu \cite{Shazeer2017}. Model \textit{embedding}
terbaru, seperti Nomic Embed \cite{NomicAI2024}, mengadopsi MoE untuk menghasilkan representasi
vektor berkualitas tinggi (seperti 768 dimensi) yang mendukung instruksi asimetris (\textit{asymmetric embedding})
untuk query dan dokumen.

% -------------------------------------------------------
\section{Metode Prototype}
% -------------------------------------------------------
Metode \textit{prototype} adalah pendekatan pengembangan perangkat lunak yang
bersifat iteratif dan interaktif, di mana sistem dikembangkan melalui siklus
penyempurnaan berdasarkan desain awal hingga tercipta versi yang mendekati
produk akhir \cite{Kholis2025}.

\begin{figure}[H]
  \centering
  \missingfig{Proses Metode Prototype (Kholis dan Zaky, 2025) -- Gambar 2.3 dari draf asli}
  \caption{Proses Metode Prototype (Sumber: Kholis \& Zaky, 2025)}
  \label{fig:prototype}
\end{figure}

% -------------------------------------------------------
\section{Black Box Testing}
% -------------------------------------------------------
\textit{Black Box Testing} adalah metode pengujian perangkat lunak yang berfokus pada
fungsionalitas sistem berdasarkan spesifikasi, tanpa memperhatikan struktur internal kode.
Tujuannya adalah memastikan bahwa input yang diberikan menghasilkan output sesuai
dengan yang diharapkan pengguna \cite{Wicaksana2025}.

% -------------------------------------------------------
\section{Basis Data}
% -------------------------------------------------------
Basis data adalah kumpulan data yang disimpan secara terstruktur dan saling berelasi,
dikelola menggunakan \textit{Database Management System} (DBMS) \cite{Santika2023}.

\subsection{Basis Data Relasional}
Basis data relasional menyimpan data dalam bentuk tabel yang saling terhubung melalui
kunci primer dan kunci asing, menggunakan SQL (\textit{Structured Query Language})
\cite{Watizega2024}. PostgreSQL merupakan salah satu RDBMS \textit{open-source} tingkat lanjut yang banyak digunakan karena ekstensibilitasnya.

\subsection{Basis Data Vektor dan pgvector}
Basis data vektor adalah sistem penyimpanan yang dioptimalkan untuk mengelola data vektor
berdimensi tinggi \cite{Wang2021}. Secara historis, solusi \textit{standalone} seperti
Milvus, Pinecone, dan Weaviate digunakan untuk tujuan ini.

Namun, \textbf{pgvector} muncul sebagai sebuah ekstensi \textit{open-source} untuk
PostgreSQL yang memungkinkan penyimpanan vektor berdimensi tinggi dan pencarian
\textit{Approximate Nearest Neighbor} (ANN) secara \textit{native} dalam basis data relasional.
Pendekatan terintegrasi ini menawarkan keunggulan arsitektural:

\begin{enumerate}[noitemsep]
  \item \textbf{Unified architecture:} data teks, vektor, dan relasional tersimpan dalam satu transaksi ACID.
  \item \textbf{Hybrid search native:} pencarian semantik dan FTS dalam satu instruksi SQL.
  \item \textbf{Kemudahan pemeliharaan:} mengurangi kompleksitas sinkronisasi antar-basis data (\textit{database drift}).
\end{enumerate}

% -------------------------------------------------------
\section{Framework dan Library Modern}
% -------------------------------------------------------
\subsection{FastAPI dan Model Asinkron}
FastAPI adalah \textit{web framework} modern untuk Python yang mendukung operasi
asinkron (\textit{async/await}) berdasarkan standar ASGI \cite{Ramirez2018}.
Keunggulan kerangka kerja ini meliputi dukungan native untuk \textit{streaming} data (seperti SSE)
dan throughput tinggi ketika diintegrasikan dengan driver basis data asinkron seperti \texttt{asyncpg},
sehingga mampu menangani ribuan koneksi secara bersamaan (\textit{high concurrency}) tanpa memblokir
\textit{thread} utama \cite{Lovering2022}.

\subsection{Next.js dan React}
Next.js adalah kerangka kerja (\textit{meta-framework}) berbasis React yang 
menyediakan fitur set tingkat produksi seperti \textit{Server-Side Rendering} (SSR), routing berbasis aplikasi (\textit{App Router}),
dan optimasi aset dinamis \cite{Adu2023}. Penggunaan kerangka kerja modern pada sisi klien (frontend)
memungkinkan pembuatan \textit{Single Page Application} (SPA) yang interaktif namun tetap memiliki performa muat yang tinggi.

\subsection{Authentication-as-a-Service}
Manajemen sesi konvensional rentan terhadap celah keamanan (seperti kebocoran kunci rahasia HS256). Pendekatan modern memanfaatkan
layanan pihak ketiga (\textit{Identity Provider}) seperti Clerk atau Auth0 yang menyediakan \textit{Authentication-as-a-Service}.
Layanan ini menangani pendaftaran pengguna dan \textit{Single Sign-On} (SSO) berbasis OAuth 2.0 \cite{Hardt2012}, kemudian menerbitkan token JWT \cite{Jones2015} menggunakan kriptografi
asimetris (RS256). Server backend (\textit{resource server}) hanya perlu memverifikasi tanda tangan JWT
tersebut menggunakan \textit{JSON Web Key Set} (JWKS) publik milik \textit{Identity Provider}.

% -------------------------------------------------------
\section{Web Crawling dan Web Scraping}
% -------------------------------------------------------
\textit{Web crawling} adalah penjelajahan sistematis halaman web oleh \textit{crawler}
untuk menemukan, mengunduh, dan mengindeks konten \cite{Khder2021}.

\textit{Web scraping} adalah proses otomatis untuk mengambil data dari situs web dan
mengubahnya dari format tidak terstruktur menjadi terstruktur \cite{Khder2021}.
Untuk menghadapi aplikasi web modern yang berbasis \textit{Single Page Application} (SPA) dan sangat
mengandalkan eksekusi JavaScript (seperti React atau Vue) untuk memuat konten, \textit{scraper} statis
seringkali gagal. Oleh karena itu, pendekatan tingkat lanjut menggunakan sistem \textit{headless browser}
(seperti Selenium WebDriver) diperlukan untuk menyimulasikan peramban nyata yang merender DOM secara
penuh sebelum proses ekstraksi dilakukan \cite{Garg2020}.

% -------------------------------------------------------
\section{Skala Likert}
% -------------------------------------------------------
Skala Likert ditemukan oleh Rensis Likert pada tahun 1932 untuk mengukur sikap,
pendapat, atau persepsi seseorang terhadap suatu objek atau pernyataan \cite{Simamora2022}.
Responden diminta menyatakan tingkat persetujuan terhadap pernyataan dalam skala
bertingkat (biasanya 1--5 atau 1--7). Skala ini secara luas diadopsi sebagai instrumen evaluasi
penerimaan pengguna (\textit{User Acceptance Testing}) dalam rekayasa perangkat lunak.

\begin{table}[H]
\centering
\caption{Skala Likert (Sumber: Simamora, 2022)}
\label{tab:likert}
\begin{tabular}{cl}
\toprule
\textbf{Skala} & \textbf{Keterangan} \\
\midrule
1 & Sangat Tidak Setuju \\
2 & Tidak Setuju \\
3 & Cukup \\
4 & Setuju \\
5 & Sangat Setuju \\
\bottomrule
\end{tabular}
\end{table}

% -------------------------------------------------------
\section{Bimbingan Karir}
% -------------------------------------------------------
Bimbingan karir adalah proses terstruktur yang membantu individu mengelola
perkembangan karir, membuat keputusan pendidikan dan pekerjaan, serta menavigasi
transisi di pasar kerja yang berubah cepat \cite{OECD2021}. UNESCO menekankan bahwa
bimbingan karir modern berperan penting dalam membekali individu dengan
\textit{career management skills} \cite{UNESCO2021}.

% -------------------------------------------------------
\section{Social Cognitive Career Theory (SCCT)}
% -------------------------------------------------------
\textit{Social Cognitive Career Theory} (SCCT) menjelaskan perkembangan minat,
pemilihan, dan pencapaian karir melalui interaksi dinamis antara keyakinan efikasi diri
(\textit{self-efficacy}), ekspektasi hasil (\textit{outcome expectations}), dan tujuan
personal \cite{Brown2023}. SCCT sangat relevan di era \textit{boundaryless career},
di mana jalur karir bersifat non-linear \cite{Wang2022}.

% -------------------------------------------------------
\section{Paradigma Life Design}
% -------------------------------------------------------
Paradigma \textit{Life Design} menekankan bahwa individu abad ke-21 secara aktif
mengonstruksi jalur karir melalui narasi yang mengintegrasikan pengalaman masa lalu,
refleksi masa kini, dan proyeksi masa depan \cite{Savickas2020}. Pendekatan ini menggeser
fokus dari sekadar mencocokkan \textit{traits} (sifat) dengan okupasi, menuju penciptaan
makna yang holistik di era kerja yang cair.

% -------------------------------------------------------
\section{Penelitian Sebelumnya}
% -------------------------------------------------------
Tabel \ref{tab:related_work} menyajikan perbandingan penelitian sebelumnya yang relevan.

\begin{longtable}{p{2.5cm}p{3cm}p{3cm}p{3cm}p{3cm}}
\caption{Penelitian Terdahulu}
\label{tab:related_work} \\
\toprule
\textbf{Peneliti} & \textbf{Judul} & \textbf{Metode/Teknologi} & \textbf{Kontribusi Utama} & \textbf{Keterbatasan} \\
\midrule
\endfirsthead
\toprule
\textbf{Peneliti} & \textbf{Judul} & \textbf{Metode/Teknologi} & \textbf{Kontribusi Utama} & \textbf{Keterbatasan} \\
\midrule
\endhead
\midrule \multicolumn{5}{r}{\textit{Lanjut di halaman berikutnya}} \\
\endfoot
\bottomrule
\endlastfoot
Nagur Vali Shaik dkk. (India, 2025) &
  \textit{Human challenging career and future advice chatbot using RAG} &
  RAG, Zephyr-7B, FAISS, LangChain, Streamlit &
  Chatbot interaktif 24/7 dengan retrieval dokumen kontekstual untuk rekomendasi karier real-time &
  Belum ada uji pengguna; tidak adaptif terhadap perubahan profil dari waktu ke waktu \\[6pt]

S. Westman (2021) &
  \textit{Artificial Intelligence for Career Guidance} &
  Kajian pustaka, tinjauan sistem AI &
  Model konseptual AI sebagai \textit{virtual coach} adaptif sepanjang hidup &
  Berbasis tinjauan teori, tanpa prototipe atau validasi pengguna \\[6pt]

RG Guntara (Indonesia, 2023) &
  \textit{Web-Based Counseling Skills Evaluation using DSRM} &
  DSR, PHP, MySQL, YouTube API &
  Sistem evaluasi keterampilan konseling berbasis DSR &
  Tidak adaptif; tanpa komponen AI canggih \\[6pt]

Haerani dkk. (2024) &
  \textit{Sistem informasi cerdas career guidance berbasis minat} &
  BFS, RMIB, PHP, MySQL &
  Rekomendasi jurusan berbasis 8 kategori minat RMIB &
  Tidak ada personalisasi adaptif; sampel UAT terbatas \\[6pt]

Elysia S. (2024) &
  \textit{Chatbot RAG untuk Layanan Pelanggan E-Commerce} &
  RAG (Llama 2 + ElasticSearch), FastAPI, React &
  Chatbot e-commerce dengan retrieval produk dan generasi naratif &
  Respons statis per sesi, tanpa pembelajaran dari histori interaksi \\
\end{longtable}

% -------------------------------------------------------
\section{Posisi Penelitian Sekarang}
% -------------------------------------------------------
Penelitian ini berfokus pada perancangan dan validasi sistem bimbingan karir adaptif
berbasis RAG dan LLM. Posisi penelitian ini ditunjukkan pada Tabel \ref{tab:position}.

\begin{table}[H]
\centering
\caption{Posisi Penelitian Ini terhadap Studi Sebelumnya}
\label{tab:position}
\begin{tabularx}{\textwidth}{Xccccccc}
\toprule
\textbf{Fitur/Kriteria} & \textbf{Shaik} & \textbf{Westman} & \textbf{Guntara} & \textbf{Haerani} & \textbf{Elysia} & \textbf{NusaNara} \\
\midrule
Arsitektur RAG \& LLM & Ya & Tidak & Tidak & Tidak & Ya & \textbf{Ya} \\
Konteks Indonesia & Tidak & Tidak & Ya & Ya & Ya & \textbf{Ya} \\
Sistem Adaptif (Dinamis) & Tidak & Ya & Tidak & Tidak & Tidak & \textbf{Ya} \\
Validasi Pengguna Formal & Tidak & Tidak & Tidak & Ya & Ya & \textbf{Ya} \\
LLM Lokal (Privasi) & Tidak & -- & -- & -- & Tidak & \textbf{Ya} \\
Hybrid Search (RRF) & Tidak & -- & -- & -- & Tidak & \textbf{Ya} \\
\bottomrule
\end{tabularx}
\end{table}

Dibandingkan penelitian sebelumnya, NusaNara memiliki keunggulan unik pada kombinasi:
arsitektur RAG dengan \textit{Hybrid Search} (RRF), LLM yang berjalan sepenuhnya lokal
(privasi data terjaga), sistem adaptif non-blocking berbasis \texttt{asyncio}, dan
basis pengetahuan yang dikurasi spesifik untuk pasar kerja Indonesia.
